% SEGMENT OLP-0489-S01
% Part: applied-modal-logics
% Chapter: epistemic-logic
% Section: public-announcement-logic-language
% SOURCE CORRECTION TA-AML-005: align the copied section metadata with this file.

\documentclass[../../../include/open-logic-section]{subfiles}

\begin{document}

\olfileid{aml}{el}{pal}

\olsection{பொது அறிவிப்புத் தருக்கம்}

காலப்போக்கிலோ புதிய தகவல் கிடைக்கும்போதோ முகவர்களின் அறிவு எவ்வாறு
மாறுகிறது என்பதைக் குறிக்க இயக்க அறிவுநிலைத் தருக்கங்கள் உதவுகின்றன.
இவற்றில் பல, தகவல்சார் \emph{நிகழ்வுகள்} அல்லது \emph{புதுப்பிப்புகள்}
வாயிலாக அறிவில் ஏற்படும் மாற்றங்களைக் குறிக்கின்றன. ஏதோவொரு வாய்பாடு
மெய்யாக அறிவிக்கப்பட்டு, அது நிகழ்வதை எல்லா முகவர்களும் ஒன்றாகக் காணும்
ஒரு பொது அறிவிப்பே மிக அடிப்படையான புதுப்பிப்பு வகை. இதைச் செய்வதற்கு,
மொழியைப் பின்வருமாறு விரிவாக்குகிறோம்.

\begin{defn}
$G$ என்பது முகவர்-குறிகளின் கணமாக இருக்கட்டும். பொது அறிவிப்புகளைக் கொண்ட
பன்முகவர் அறிவுநிலைத் தருக்கத்தின் அடிப்படை மொழி பின்வருவனவற்றைக்
கொண்டுள்ளது:
\begin{enumerate}
  \tagitem{prvFalse}{!!{falsity}க்கான கூற்று மாறிலி~$\lfalse$.}{}
  \tagitem{prvTrue}{!!{truth}க்கான கூற்று மாறிலி~$\ltrue$.}{}
  \item !!{propositional variable}s இன் ஓர் !!{denumerable}s கணம்: $\Obj
    p_0$, $\Obj p_1$, $\Obj p_2$, \dots
  \item முன்மொழிவு இணைப்பிகள்: \startycommalist
  \iftag{prvNot}{\ycomma $\lnot$ (மறுப்பு)}{}%
  \iftag{prvAnd}{\ycomma $\land$ (இணைப்பு)}{}%
  \iftag{prvOr}{\ycomma $\lor$ (பிரிப்பு)}{}%
  \iftag{prvIf}{\ycomma $\lif$ (!!{conditional})}{}%
  % SOURCE CORRECTION TA-AML-006: the formation clauses include biconditionals.
  \iftag{prvIff}{\ycomma $\liff$ (!!{biconditional})}{}.
  \item $a \in G$ ஆகும் இடத்தில் அறிவுச் செயற்குறி~$\Knows_a$.
  \item $!B$ ஓர் !!a{formula} ஆகும் இடத்தில் பொது அறிவிப்புச்
    செயற்குறி~$[!B]$.
\end{enumerate}
\end{defn}

பொது அறிவிப்புச் செயற்குறி ஒரு பெட்டிச் செயற்குறியாகச் செயல்படுகிறது;
அதற்கேற்ப மொழியின் தொகுமுறை வரையறை பின்வருமாறு அளிக்கப்படுகிறது:

\begin{defn}
  அறிவுநிலை மொழியின் \emph{!!^{formula}s} பின்வருமாறு தொகுமுறையில்
  வரையறுக்கப்படுகின்றன:
  \begin{enumerate}
  \tagitem{prvFalse}{$\lfalse$ ஓர் அணு !!{formula}.}{}

  \tagitem{prvTrue}{$\ltrue$ ஓர் அணு !!{formula}.}{}

  \item ஒவ்வொரு கூற்று மாறி $\Obj p_i$ உம் ஓர் (அணு)
  !!{formula}.

  \tagitem{prvNot}{$!A$ ஓர் !!{formula} எனில், $\lnot !A$ உம் ஓர்
  !!{formula}.}{}

  \tagitem{prvAnd}{$!A$, $!B$ ஆகியவை !!{formula}s எனில், $(!A \land
  !B)$ ஓர் !!{formula}.}{}

  \tagitem{prvOr}{$!A$, $!B$ ஆகியவை !!{formula}s எனில், $(!A \lor !B)$
  ஓர் !!{formula}.}{}

  \tagitem{prvIf}{$!A$, $!B$ ஆகியவை !!{formula}s எனில், $(!A \lif !B)$
  ஓர் !!{formula}.}{}

  \tagitem{prvIff}{$!A$, $!B$ ஆகியவை !!{formula}s எனில், $(!A \liff !B)$
  ஓர் !!{formula}.}{}

  \item $!A$ ஓர் !!{formula} ஆகவும் $a \in G$ ஆகவும் இருந்தால்,
  $\Knows_a !A$ ஓர் !!{formula}.

  \item $!A$, $!B$ ஆகியவை !!{formula}s எனில், $[!A] !B$ ஓர் !!a{formula}.

  \tagitem{limitClause}{வேறு எதுவும் ஓர் !!a{formula} அல்ல.}{}
\end{enumerate}
\end{defn}

!!{formula}~$[!A] !B$ இன் நோக்கப்பட்ட வாசிப்பு, ``$!A$ மெய்யாக
அறிவிக்கப்பட்ட பிறகு, $!B$~மெய்யாகிறது'' என்பதாகும். பொது அறிவிப்புகளின்
சூழலில் பொதுவறிவைப் பற்றிப் பேசுவதும் சில நேரங்களில் பயனுள்ளதாக இருக்கும்;
எனவே பொதுவறிவுச் செயற்குறி~$\CKnows_G !A$ ஐயும் மொழி கொண்டிருக்கலாம்.

\end{document}
